\section{Related Work} \label{sec:related_work}

Plagiarism detection has evolved significantly with advances in machine learning and natural language processing. Early systems relied on string-matching and hash-based techniques, which are effective for detecting near-duplicate content but vulnerable to paraphrasing attacks.

\subsection{Traditional Approaches}

Fingerprinting-based systems such as MOSS (Measure of Software Similarity) and Turnitin employ hash-based matching, which excels at detecting direct copying with $O(n)$ complexity. However, these approaches inherently fail against lexical variation and semantic paraphrasing, as demonstrated in \cite{Schleimer2003} with the Rabin fingerprinting algorithm.

String-matching algorithms including Longest Common Subsequence (LCS) and Levenshtein distance provide character-level similarity but lack semantic understanding. Cosine similarity with bag-of-words representations marked a transition toward semantic awareness, though limited by vocabulary sparsity.

\subsection{Machine Learning Era}

TF-IDF (Term Frequency-Inverse Document Frequency) combined with cosine similarity became a standard baseline, offering improved efficiency over string matching. \cite{Robertson2009} established TF-IDF's strong performance for document similarity tasks. Nevertheless, these methods ignore word order, syntactic structure, and semantic relationships.

Support Vector Machines (SVMs) and Random Forests were applied to plagiarism detection with engineered features (n-grams, stylometric features), as documented in \cite{Stamatatos2009}. These approaches improved upon TF-IDF but remained constrained by hand-crafted feature engineering.

\subsection{Deep Learning Approaches}

Convolutional Neural Networks (CNNs) and Recurrent Neural Networks (RNNs) introduced learned representations, with BiLSTM architectures capturing sequential dependencies. \cite{Hochreiter1997} established LSTM's superiority for long-range dependencies, which BiLSTM extends with bidirectional context.

Word embeddings including Word2Vec \cite{Mikolov2013} and GloVe \cite{Pennington2014} further improved semantic representation. These embeddings capture word relationships, enabling models to detect paraphrased plagiarism through semantic similarity in embedding space.

\subsection{Transformer-based Models}

BERT (Bidirectional Encoder Representations from Transformers) \cite{Devlin2018} revolutionized NLP through bidirectional pre-training on massive text corpora. The self-attention mechanism \cite{Vaswani2017} enables capturing long-range dependencies and complex semantic relationships superior to RNNs.

Recent work applies BERT to plagiarism detection with fine-tuning for binary classification. \cite{Shivakumar2020} demonstrated BERT's effectiveness for source code plagiarism detection. \cite{Hosseini2021} extended this to natural language plagiarism with multilingual variants.

\subsection{Gap Analysis}

While individual approaches are well-documented, systematic comparative evaluation of classical ML, deep learning, and transformer methods remains limited. Most research focuses on single-method evaluation without comprehensive efficiency analysis. Our work addresses this gap by providing detailed comparative analysis across three distinct paradigms with computational trade-off quantification.
